\documentclass[../../main.tex]{subfiles}

\begin{document}

\chapter{Kriptografi Kunci-Publik: Algoritma RSA}
\label{chap:rsa}

\begin{learningoutcome}
Setelah mempelajari bab ini, mahasiswa diharapkan mampu:
\begin{itemize}
    \item Menjelaskan pengertian kriptografi kunci-publik dan penggunaan sepasang kunci pada algoritma RSA (Sub-CPMK 1.3).
    \item Menjelaskan prosedur pembangkitan pasangan kunci publik dan privat pada algoritma RSA (Sub-CPMK 1.3).
    \item Menjelaskan rumus enkripsi dan dekripsi RSA untuk mengamankan pesan (Sub-CPMK 1.3).
    \item Menjelaskan landasan matematis RSA, termasuk Teorema Euler dan Fungsi Totient Euler (Sub-CPMK 1.3).
    \item Menjelaskan keamanan RSA berdasarkan sulitnya pemfaktoran bilangan bulat besar (Sub-CPMK 1.3).
    \item Mengidentifikasi risiko serangan \textit{Man-in-the-Middle} (MITM) dalam distribusi kunci publik (Sub-CPMK 1.3).
\end{itemize}
\end{learningoutcome}

\subfile{section-01}
\subfile{section-02}
\subfile{section-03}
\subfile{section-04}
\section{Contoh Terhitung}
\label{sec:contoh-rsa}

Contoh-contoh berikut memakai parameter yang sama, yaitu $p = 47$, $q = 71$,
dan $e = 79$. Kunci publiknya $(e, n) = (79, 3337)$ dan kunci privatnya
$(d, n) = (1019, 3337)$. Modulusnya sengaja dipilih kecil supaya setiap
langkah dapat diperiksa dengan tangan.

\begin{example}[Pembangkitan kunci RSA]
\label{ex:rsa-kunci}
Alice memilih dua bilangan prima kecil $p = 47$ dan $q = 71$.
\begin{enumerate}
    \item Hitung modulus:
    \[ n = p \cdot q = 47 \times 71 = 3337. \]
    \item Hitung fungsi Totient Euler:
    % Koreksi: slide menulis 46 x 60 = 3220. Hasil yang benar adalah
    % 46 x 70 = 3220.
    \[ \varphi(n) = (p-1)(q-1) = 46 \times 70 = 3220. \]
    \item Pilih kunci publik $e = 79$. Syaratnya terpenuhi karena
          $\text{PBB}(79, 3220) = 1$.
    \item Hitung kunci privat $d$ sebagai invers $e$ modulo $\varphi(n)$:
    \[ e \cdot d \equiv 1 \pmod{3220}
       \;\Longrightarrow\; d = \frac{1 + k\,\varphi(n)}{e}
       = \frac{1 + 3220k}{79}. \]
    Nilai $k = 1, 2, 3, \dots$ dicoba sampai pembilangnya habis dibagi 79.
    Untuk $k = 1$ hasilnya $3221/79 = 40{,}77$ dan untuk $k = 2$ hasilnya
    $6441/79 = 81{,}53$. Keduanya bukan bilangan bulat. Untuk $k = 25$,
    \[ d = \frac{1 + 3220 \times 25}{79} = \frac{80501}{79} = 1019. \]
    Pemeriksaan: $79 \times 1019 = 80501 = 1 + 25 \times 3220$.
\end{enumerate}
Kunci publik Alice adalah $(e, n) = (79, 3337)$ dan kunci privatnya
$(d, n) = (1019, 3337)$. Hanya $n$ dan $e$ yang boleh disebarkan.
\end{example}

\begin{example}[Enkripsi dan dekripsi RSA]
Bob mengirim pesan $m = 16$ kepada Alice dengan kunci publik $(79, 3337)$.
\begin{enumerate}
    \item Enkripsi oleh Bob:
    \[ c = m^e \bmod n = 16^{79} \bmod 3337 = 1493. \]
    \item Dekripsi oleh Alice dengan kunci privat $(1019, 3337)$:
    \[ m = c^d \bmod n = 1493^{1019} \bmod 3337 = 16. \]
\end{enumerate}
Plainteks $m = 16$ pulih dengan tepat. Bilangan 16 di sini adalah bilangan
bulat, bukan huruf. Cara mengubah teks menjadi bilangan bulat dibahas pada
contoh berikutnya.
\end{example}

\begin{example}[Enkripsi pesan multi-blok ``HELLO ALICE'']
\label{ex:rsa-hello-alice}
Bob ingin mengirim pesan $M = \text{HELLO ALICE}$ kepada Alice. Pesan itu
diubah lebih dahulu menjadi bilangan bulat, lalu dipecah menjadi blok-blok
yang lebih kecil daripada modulus.

\textbf{Langkah 1: pengodean.} Setiap huruf dikodekan menjadi dua digit dengan
aturan $\text{A} = 00$, $\text{B} = 01$, sampai $\text{Z} = 25$. Spasi
diabaikan. Tabel~\ref{tab:pengodean-hello} memperlihatkan hasilnya.

\begin{table}[htbp]
    \centering
    \begin{tabular}{@{}l *{10}{c}@{}}
    \toprule
    Huruf & H & E & L & L & O & A & L & I & C & E \\
    \midrule
    Kode & 07 & 04 & 11 & 11 & 14 & 00 & 11 & 08 & 02 & 04 \\
    \bottomrule
    \end{tabular}
    \caption{Pengodean pesan ``HELLO ALICE'' dengan aturan $\text{A}=00$ sampai
    $\text{Z}=25$. Spasi tidak dikodekan}
    \label{tab:pengodean-hello}
\end{table}

% Koreksi: slide menulis 070411111140011080204 (21 digit). Pengodean yang
% benar menghasilkan 20 digit, yaitu 07041111140011080204.
Deretan kodenya menjadi $m = 07041111140011080204$, yaitu bilangan 20 digit.
Nilai itu lebih besar daripada $n = 3337$, sehingga harus dipecah. Dengan
mengambil empat digit per blok diperoleh $m_1 = 0704$, $m_2 = 1111$,
$m_3 = 1400$, $m_4 = 1108$, dan $m_5 = 0204$. Kelimanya memenuhi syarat
$0 \le m_i \le 3336$ pada Persamaan~\eqref{eq:syarat-blok}.

\textbf{Langkah 2: enkripsi.} Setiap blok dienkripsi dengan kunci publik Alice,
$c_i = m_i^{79} \bmod 3337$. Hasilnya adalah
% Koreksi: slide menulis m_1 = 204 pada baris kelima. Yang dimaksud adalah
% m_5 = 0204; hasil c_5 = 1164 sudah benar.
\[
\begin{aligned}
m_1 = 0704 &\quad\longrightarrow\quad c_1 = 704^{79} \bmod 3337 = 0328, \\
m_2 = 1111 &\quad\longrightarrow\quad c_2 = 1111^{79} \bmod 3337 = 0301, \\
m_3 = 1400 &\quad\longrightarrow\quad c_3 = 1400^{79} \bmod 3337 = 2653, \\
m_4 = 1108 &\quad\longrightarrow\quad c_4 = 1108^{79} \bmod 3337 = 2986, \\
m_5 = 0204 &\quad\longrightarrow\quad c_5 = 204^{79} \bmod 3337 = 1164.
\end{aligned}
\]
Bob mengirimkan cipherteks $C = 0328, 0301, 2653, 2986, 1164$ kepada Alice.
Setiap blok ditulis empat digit. Nol di depan pada $0704$, $0328$, $0301$, dan
$0204$ tetap ditulis agar seluruh blok sejajar.

\textbf{Langkah 3: dekripsi.} Alice memakai kunci privatnya,
$m_i = c_i^{1019} \bmod 3337$, dan memperoleh kembali kelima blok semula.
Tabel~\ref{tab:jejak-hello-alice} merangkum perjalanan pesan dari huruf sampai
huruf kembali.

\begin{table}[htbp]
    \centering
    \begin{tabular}{@{}c c c c l@{}}
    \toprule
    \textbf{Blok} & \textbf{Plainteks $m_i$} & \textbf{Cipherteks $c_i$} & \textbf{Pulih} &
    \textbf{Huruf} \\
    \midrule
    1 & \code{0704} & \code{0328} & \code{0704} & H, E \\
    2 & \code{1111} & \code{0301} & \code{1111} & L, L \\
    3 & \code{1400} & \code{2653} & \code{1400} & O, A \\
    4 & \code{1108} & \code{2986} & \code{1108} & L, I \\
    5 & \code{0204} & \code{1164} & \code{0204} & C, E \\
    \bottomrule
    \end{tabular}
    \caption{Jejak enkripsi dan dekripsi pesan ``HELLO ALICE'' dengan kunci
    publik $(79, 3337)$ dan kunci privat $(1019, 3337)$}
    \label{tab:jejak-hello-alice}
\end{table}

Deretan angka yang dipulihkan adalah $07041111140011080204$. Deretan itu
diterjemahkan kembali menjadi $\text{HELLO ALICE}$. Spasi tidak dikodekan,
sehingga penerima menambahkannya kembali.
\end{example}


\begin{obeactivity}
\textbf{Aktivitas 10.1: Simulasi RSA dengan Bilangan Prima Kecil}
1. Pilih dua bilangan prima kecil, misalnya $p = 3$ dan $q = 11$.
2. Hitung $n = pq$ dan $\varphi(n) = (p-1)(q-1)$.
3. Pilih $e = 3$. Periksa bahwa $\text{PBB}(e, \varphi(n)) = 1$.
4. Hitung $d$ dengan mencoba nilai $k = 1, 2, 3, \dots$ pada
   $d = \bigl(1 + k\varphi(n)\bigr)/e$ sampai hasilnya bilangan bulat.
5. Enkripsi pesan $m = 5$ dengan $(e, n)$, lalu dekripsi hasilnya dengan
   $(d, n)$. Periksa bahwa pesan pulih dengan tepat.
6. Ulangi langkah 3 dengan $e = 5$. Jelaskan mengapa langkah 4 tidak dapat
   dijalankan pada keadaan itu.
\end{obeactivity}

Aktivitas 10.1 dikerjakan dengan tangan. Pada Aktivitas 10.2, langkah
pengodean dan pemecahan blok dapat diperiksa dengan program pada
Listing~\ref{lst:rsa-pesan}.

\begin{obeactivity}
\textbf{Aktivitas 10.2: Enkripsi dan Dekripsi Pesan Multi-Blok}
1. Kodekan nama Anda sendiri, paling sedikit sepuluh huruf, dengan aturan
   $\text{A} = 00$ sampai $\text{Z} = 25$.
2. Pecah deretan kode itu menjadi blok empat digit. Periksa setiap blok
   memenuhi syarat pada Persamaan~\eqref{eq:syarat-blok} untuk $n = 3337$.
3. Enkripsi setiap blok dengan kunci publik $(79, 3337)$. Tuliskan cipherteks
   dalam empat digit, termasuk nol di depan bila ada.
4. Dekripsi kembali cipherteks Anda dengan kunci privat $(1019, 3337)$, lalu
   terjemahkan hasilnya menjadi huruf. Apakah nama Anda pulih dengan tepat?
5. Hapus nol di depan pada salah satu blok cipherteks, lalu dekripsi deretan
   angka yang sudah berubah itu. Jelaskan mengapa hasilnya tidak lagi sama.
\end{obeactivity}

\lstinputlisting[
    language=Python,
    caption={Pengodean pesan, pemecahan blok, dan penjagaan angka nol di depan},
    label={lst:rsa-pesan},
    firstline=73,
    lastline=123
]{\codefile{python/rsa_uji.py}}

\begin{obeactivity}
\textbf{Aktivitas 10.3: Menimbang Ukuran Modulus}
1. Salin pernyataan bahan kuliah tentang modulus 256 bit, 512 bit, dan
   1024 bit.
2. Jelaskan mengapa panjang modulus mempengaruhi usaha pemfaktoran.
3. Pada tahun 1977, ketiga penemu RSA menantang pembaca \textit{Scientific
   American} untuk memecahkan cipherteks RSA. Tantangan itu baru terpecahkan
   pada tahun 1994. Diskusikan apa yang ditunjukkan jeda waktu tersebut
   tentang perkiraan keamanan RSA.
4. Jelaskan mengapa algoritma Shor mengubah ancaman itu, meskipun panjang
   modulus tetap diperbesar.
\end{obeactivity}


\begin{obereflection}
\textbf{Latihan 10.1}
1. Sebutkan tiga penemu algoritma RSA, lembaganya, dan tahun pengenalannya.
2. Jelaskan mengapa keamanan RSA terletak pada sulitnya memfaktorkan $n$
   menjadi $p$ dan $q$.
3. Dengan $p = 47$ dan $q = 71$, hitunglah $\varphi(n)$. Periksa bahwa
   $e = 17$ memenuhi syarat sebagai kunci publik, lalu tentukan $d$ yang
   bersesuaian.
4. Jelaskan mengapa $p$, $q$, dan $\varphi(n)$ harus dirahasiakan, sedangkan
   $n$ boleh diketahui umum.
5. Dengan kunci publik $(79, 3337)$ dan kunci privat $(1019, 3337)$,
   enkripsikan $m = 16$, lalu dekripsikan hasilnya.
6. Mengapa pesan yang lebih panjang daripada modulus harus dipecah menjadi
   blok dengan syarat $0 \le m_i \le n-1$?
\end{obereflection}

\begin{obereflection}
\textbf{Latihan 10.2}
1. Jelaskan mengapa RSA lebih lambat daripada algoritma kunci-simetri, lalu
   jelaskan cara kriptografi hibrida mengatasi hal itu.
2. Uraikan jalannya serangan \textit{Man-in-the-Middle} ketika Alice dan Bob
   bertukar kunci publik.
3. Jelaskan mengapa \textit{padding} OAEP membuat cipherteks menjadi tidak
   deterministik.
4. Bahan kuliah menyatakan bahwa modulus 256 bit dapat difaktorkan dalam
   beberapa jam dan modulus 512 bit dengan beberapa ratus komputer. Apa
   implikasi pernyataan itu bagi pemilihan panjang modulus RSA?
5. Pada tahun 1977 ketiga penemu RSA membuat sayembara di \textit{Scientific
   American}. Kapan sayembara itu terpecahkan, dan berapa lama waktu yang
   dibutuhkan?
6. Jelaskan mengapa algoritma Shor mengancam RSA, dan mengapa memperbesar
   modulus tidak menangkal ancaman itu.
\end{obereflection}

\section{Rangkuman}
\begin{itemize}
    \item RSA adalah algoritma kunci-publik yang paling terkenal dan paling
          banyak dipakai. Kunci publiknya $(e, n)$ dan kunci privatnya
          $(d, n)$.
    \item $p$ dan $q$ adalah bilangan prima rahasia, $n = pq$ boleh diketahui
          umum, $\varphi(n) = (p-1)(q-1)$ harus rahasia, $e$ adalah kunci
          publik dengan $\text{PBB}(e, \varphi(n)) = 1$, dan $d$ adalah kunci
          privat dengan $ed \equiv 1 \pmod{\varphi(n)}$.
    \item Keamanan RSA terletak pada sulitnya memfaktorkan $n$ menjadi $p$ dan
          $q$. Setelah $n$ difaktorkan, $\varphi(n)$ dan $d$ dapat dihitung.
    \item Fungsi Totient Euler untuk $n = pq$ adalah $\varphi(n) = (p-1)(q-1)$.
          Teorema Euler, $a^{\varphi(n)} \equiv 1 \pmod n$ bila
          $\text{PBB}(a,n)=1$, dipakai untuk menurunkan rumus RSA.
    \item Penurunannya berjalan dari $e \cdot d = k\varphi(n) + 1$ menuju
          $m^{ed} \equiv m \pmod n$. Jadi enkripsi dan dekripsi saling
          membalik.
    \item Pembangkitan kunci: pilih $p$ dan $q$ prima, hitung $n = pq$ dan
          $\varphi(n)$, pilih $e$ yang relatif prima terhadap $\varphi(n)$,
          lalu hitung $d \equiv e^{-1} \pmod{\varphi(n)}$. Nilai $e = 65537$
          lazim dipakai karena dianggap aman dan efisien.
    \item Enkripsi memakai $c = m^e \bmod n$ dan dekripsi memakai
          $m = c^d \bmod n$. Pesan yang lebih panjang daripada modulus dipecah
          menjadi blok dengan syarat $0 \le m_i \le n-1$.
    \item Setiap blok ditulis dengan panjang tetap. Nol di depan, seperti pada
          $0704$ dan $0204$, tetap ditulis agar batas antrablok tidak
          bergeser.
    \item Contoh: $p=47$, $q=71$ memberi $n=3337$, $\varphi(n)=3220$,
          $e=79$, dan $d=1019$. Pesan ``HELLO ALICE'' dikodekan menjadi
          $07041111140011080204$, dipecah menjadi lima blok, dienkripsi
          menjadi $0328, 0301, 2653, 2986, 1164$, lalu pulih dengan tepat.
    \item Modulus 1024 bit setara 300 angka desimal. Modulus 256 bit dapat
          difaktorkan dalam beberapa jam, sedangkan modulus 512 bit memerlukan
          beberapa ratus komputer. Bahan kuliah ini menyatakan bahwa $n$ di
          atas 1024 bit baru dianggap aman.
    \item Sayembara RSA pada tahun 1977 di majalah \textit{Scientific American}
          baru terpecahkan pada tahun 1994 dalam waktu 8 bulan.
    \item Komputasi kuantum mengancam RSA karena algoritma Shor dapat
          memfaktorkan $n$ dalam waktu polinomial.
    \item RSA lebih lambat daripada algoritma kunci-simetri. Dalam praktik,
          RSA dipakai untuk mengenkripsi kunci sesi, sedangkan pesan
          dienkripsi dengan algoritma simetri. Kombinasi itu dinamakan
          kriptografi hibrida.
    \item OAEP adalah metode \textit{padding} yang membuat pesan menjadi acak
          sebelum dienkripsi. Cipherteksnya tidak deterministik dan aman
          terhadap \textit{chosen-ciphertext attack}.
    \item Distribusi kunci publik dapat diserang \textit{Man-in-the-Middle}.
          Carol dapat menyamar sebagai Alice dan Bob, lalu membaca pesan tanpa
          diketahui kedua pihak.
\end{itemize}


\begin{obeassessment}
\textbf{Kuis Bab 10}
1. Jelaskan peran fungsi Totient Euler dalam pembangkitan kunci RSA.
2. Bagaimana Teorema Euler dipakai untuk menunjukkan bahwa dekripsi RSA
   mengembalikan pesan asli?
3. Dengan $p = 5$, $q = 11$, dan $e = 3$, hitunglah $n$, $\varphi(n)$, dan
   $d$. Kemudian enkripsikan $m = 9$ dan periksa bahwa dekripsinya memulihkan
   $m$.
4. Jelaskan mengapa modulus RSA harus panjang. Sebutkan keadaan modulus 256
   bit, 512 bit, dan 1024 bit menurut bahan kuliah.
5. Jelaskan mengapa RSA jarang dipakai untuk mengenkripsi data berukuran
   besar, dan bagaimana kriptografi hibrida mengatasi hal itu.
6. Mengapa RSA memerlukan \textit{padding} seperti OAEP sebelum pesan
   dienkripsi?
\end{obeassessment}
\begin{competencychecklist}
\begin{tabularx}{\linewidth}{|X|c|}
\hline
Indikator Kompetensi & Tercapai \\
\hline
Saya dapat membedakan kunci publik dan kunci privat pada RSA & [ ] \\
\hline
Saya dapat menjelaskan properti dan kerahasiaan komponen algoritma RSA & [ ] \\
\hline
Saya dapat menjelaskan Teorema Euler dan Fungsi Totient Euler & [ ] \\
\hline
Saya dapat menghitung pasangan kunci RSA untuk parameter kecil & [ ] \\
\hline
Saya dapat melakukan proses enkripsi dan dekripsi RSA & [ ] \\
\hline
Saya dapat mengenkripsi pesan multi-blok dengan menjaga angka nol di depan & [ ] \\
\hline
Saya dapat menjelaskan kaitan pemfaktoran bilangan bulat dengan keamanan RSA & [ ] \\
\hline
Saya dapat menjelaskan pengaruh panjang modulus dan komputasi kuantum & [ ] \\
\hline
Saya dapat menjelaskan mekanisme kriptografi hibrida & [ ] \\
\hline
Saya dapat menjelaskan peran \textit{padding} OAEP & [ ] \\
\hline
Saya dapat menjelaskan serangan \textit{Man-in-the-Middle} pada distribusi kunci publik & [ ] \\
\hline
\end{tabularx}
\end{competencychecklist}


\end{document}
